\documentclass[10pt,a4paper]{article}
\usepackage[T1]{fontenc}
\usepackage[margin=2cm]{geometry}
\usepackage{isabelle,isabellesym}
\usepackage{amsmath,amssymb}
\usepackage{pdfsetup}

\urlstyle{rm}
\isabellestyle{it}

\begin{document}

\title{Quantum Speed Limits: the Zieli\'nski--Zych Bound and Time--Occupation Bounds for Qubit--Bus--Qubit State Transfer}
\author{Carles Mar\'in}
\maketitle

\begin{abstract}
We prove the family of Zieli\'nski--Zych speed limits~\cite{ZZ06}, $\pi^a\le 2M_a\tau^a$ for every exponent
$0<a\le 1$, where $M_a$ is the fractional energy moment and $\tau$ a time at which the autocorrelation vanishes,
together with its limit $a\to 0$. For the relay qubit--bus--qubit with couplings bounded by $|g|\le G$, we prove that
the transfer time satisfies $GT\ge\pi/\sqrt2$, by means of an explicit calibration function; we prove the dephasing
excess $\tfrac23 q(1-q)$ over the additive budget and the exact integrals of the optimal protocol, $3\pi/8$,
$\pi/4$ and their corrected values; and we prove the lower bound $(1-\varepsilon)/(G^2T)$ on the integrated
occupation of the bus for any control, and the time--occupation frontier $2u\le G^2T\int p$. With a lossy bus
(excess loss rate $\delta$) we prove the exact error floor $P(T)\le e^{-\gamma T}/(1+\delta/(G^2T))$ for any
control and any duration, and the time--occupation frontier for partial transfer, whose constant tends to~2. The trajectories
are continuous and differentiable outside a countable set of times.
\end{abstract}

\paragraph{Acknowledgement.} The formalization and the text of this document were prepared with the
assistance of generative AI systems, and checked by the author.

\tableofcontents

\section*{Introduction}

The theory \texttt{Zielinski\_Zych} proves the speed limits of Zieli\'nski and Zyczkowski~\cite{ZZ06} for a discrete
spectral measure, and the behaviour of the bound as the exponent tends to zero. The remaining theories concern a relay:
a qubit~A is coupled to a bus mode and the mode to a qubit~B, with coupling moduli at most~$G$. Let $x=|a_A|$ and
$z=|a_B|$ be the moduli of the amplitudes of the two qubits; then $1-x^2-z^2$ is the occupation of the bus.
\texttt{Relay\_Minimal\_Time} gives the minimal transfer time, \texttt{Relay\_Dephasing} the dephasing excess and
the exact integrals of the optimal protocol, and \texttt{Relay\_Occupation} the lower bound on the integrated
occupation and the frontier. The theory \texttt{Relay\_Loss\_Floor} adds a lossy bus: the exact error floor
$P(T)\le e^{-\gamma T}/(1+\delta/(G^2T))$ and the time--occupation frontier for partial transfer. For time-optimal control of three-level systems see~\cite{BCC05}; for STIRAP-based
state transfer through a bus under dissipation see~\cite{WYC17}; for the experimental context
see~\cite{SUN26,MAR26}.

\parindent 0pt\parskip 0.5ex

The theory \texttt{Relay\_Endpoints} proves that the end qubits must be occupied: $G$ times the
integrated occupation of emitter and receiver is at least $\pi/2$ for a complete transfer, for
absolutely continuous amplitudes; \texttt{Relay\_Windows} gives the comparison windows used when the
qubits decay faster than the bus.

\input{session}

\bibliographystyle{abbrv}
\bibliography{root}

\end{document}
